\documentclass{article}
\usepackage[T1]{fontenc}
\usepackage[utf8]{inputenc}
\usepackage{polski}
\usepackage{listings}
\usepackage{pgfplots}
\usepackage{float}
\usepackage{booktabs}
\usepackage{pgfplotstable}
\usepackage{xcolor}
\usepackage{geometry}
\usepackage{amsmath,amssymb}
\usepackage{caption}
\usepgfplotslibrary{fillbetween}
\geometry{margin=2.2cm}
\pgfplotsset{compat=1.18}

% --- Definicje kolorów ---
\definecolor{cBlue}{RGB}{33,113,181}
\definecolor{cOrange}{RGB}{230,85,13}
\definecolor{cGreen}{RGB}{49,163,84}
\definecolor{cPurple}{RGB}{117,107,177}
\definecolor{cBrown}{RGB}{166,97,26}
\definecolor{cPink}{RGB}{231,41,138}

% --- Styl wykresów ---
\pgfplotsset{
	myaxis/.style={
		width=0.9\linewidth, height=7cm,
		xlabel={$n$ (rozmiar danych)}, ylabel={czas [$\mu$s/ms]},
		grid=both,
		legend pos=north west, legend cell align=left,
		title style={yshift=6pt}
	},
	linearxy/.style={xmode=linear, ymode=linear},
	logxy/.style={xmode=log, ymode=log}
}

% --- Styl kodu ---
\lstset{
	language=C++,
	basicstyle=\ttfamily\small,
	keywordstyle=\bfseries,
	commentstyle=\itshape,
	numbers=left,
	numberstyle=\tiny,
	stepnumber=1,
	numbersep=8pt,
	showstringspaces=false,
	breaklines=true,
	frame=single,
	tabsize=2
}

\title{Analiza algorytmów: Programowanie Dynamiczne i Zachłanne\\(Rod Cutting, LCS, Activity Selector, Huffman)}
\author{}
\date{}

\begin{document}
	\maketitle
	
	\section{Wprowadzenie}
	W sprawozdaniu przedstawiono implementację i analizę wydajności algorytmów rozwiązujących cztery klasyczne problemy optymalizacyjne w zależności od różnych podejść (naiwnych, zachłannych, dynamicznych).
	
	\section{Metodyka pomiarów}
	Pomiary czasu wykonania zaimplementowano w języku C++ przy użyciu biblioteki \texttt{chrono}. Dla każdego algorytmu przeprowadzono serię testów na losowych danych wejściowych o różnym rozmiarze $n$. Wyniki uśredniono, aby zminimalizować błąd pomiarowy.
	
	\section{Implementacja i Analiza}
	
	% ==========================================
	% 1. ROD CUTTING
	% ==========================================
	\subsection{Cięcie pręta (Rod Cutting)}
	Problem polega na znalezieniu sposobu pocięcia pręta o długości $n$, aby zmaksymalizować zysk, mając dany cennik.
	
	\subsubsection{Implementacja}
	Zaimplementowano trzy wersje algorytmu:
	\begin{enumerate}
		\item \textbf{Naiwna (rekurencyjna):} $O(2^n)$ – sprawdza wszystkie możliwe podziały.
		\item \textbf{Memoized (rekurencja ze spamiętywaniem):} $O(n^2)$ – zapamiętuje wyniki podproblemów.
		\item \textbf{Iteracyjna (Bottom-Up):} $O(n^2)$ – wypełnia tablicę od najmniejszych podproblemów.
	\end{enumerate}
	\newpage
	
	\begin{lstlisting}[caption={Fragment implementacji: Rod Cutting (wersja iteracyjna)}]
		// Algorytm Bottom-Up o zlozonosci O(n^2).
		// Rozwiazuje podproblemy od najmniejszej dlugosci preta (j) do n.
		Solution cut_rod_bottom_up(const vector<int>& p, int n) {
			vector<int> revenue(n + 1); // Tablica zyskow 
			vector<int> sol(n + 1); // Tablica rekonstrukcji rozwiazania
			
			revenue[0] = 0;
			
			// Iterujemy po dlugosciach preta od 1 do n
			for (int j = 1; j <= n; ++j) {
				int q = INT_MIN;
				// Sprawdzamy kazde mozliwe pierwsze ciecie 'i' (od 1 do j)
				for (int i = 1; i <= j; ++i) {
					// p[i-1] to cena kawalka o dlugosci i
					// revenue[j-i] to optymalny zysk dla pozostalej czesci
					int val = p[i - 1] + revenue[j - i];
					
					if (q < val) {
						q = val;
						sol[j] = i; // Zapamietujemy optymalne ciecie dla dlugosci j
					}
				}
				revenue[j] = q;
			}
			
			return {revenue[n], sol};
		}
	\end{lstlisting}
	
	\subsubsection{Wyniki pomiarów}
	Poniższa tabela oraz wykres przedstawiają czasy wykonania (w mikrosekundach) dla różnych długości pręta.
	
	% Tabela Rod Cutting - POPRAWIONE NAZWY KOLUMN
	\begin{table}[H]
		\centering
		\caption{Wyniki pomiarów dla algorytmów Rod Cutting}
		\pgfplotstabletypeset[
		col sep=comma,
		string type,
		% Uzywamy nazw kolumn z CSV: N, Naiwny, Memoized, Iteracyjny
		columns={N,Naiwny,Memoized,Iteracyjny},
		columns/N/.style={column name=$n$, column type={r}},
		columns/Naiwny/.style={column name={Naiwny [$\mu$s]}, column type={r}},
		columns/Memoized/.style={column name={Memoized [$\mu$s]}, column type={r}},
		columns/Iteracyjny/.style={column name={Iteracyjny [$\mu$s]}, column type={r}},
		every head row/.style={before row=\toprule,after row=\midrule},
		every last row/.style={after row=\bottomrule},
		]{rod.csv}
	\end{table}
	
	\begin{figure}[H]
		\centering
		\begin{tikzpicture}
			\begin{axis}[myaxis, title={Rod Cutting: Porównanie rozwiązań dynamicznych}, xlabel={$n$}, ylabel={czas [$\mu$s]}]
				
				\addplot+[line width=1pt, color=cGreen, mark=triangle*, mark size=2pt]
				table[x=N, y=Naiwny, col sep=comma, restrict x to domain=0:26]{rod.csv};
				\addlegendentry{Naiwny}
				
				% Wykres: Memoized
				\addplot+[line width=1pt, color=cBlue, mark=*, mark size=2pt]
				table[x=N, y=Memoized, col sep=comma]{rod.csv};
				\addlegendentry{Memoized}
				
				% Wykres: Iteracyjny
				\addplot+[line width=1pt, color=cOrange, mark=square*, mark size=2pt]
				table[x=N, y=Iteracyjny, col sep=comma]{rod.csv};
				\addlegendentry{Iteracyjny}
				
			\end{axis}
		\end{tikzpicture}
		\caption{Dla dużych $n$ obie wersje dynamiczne działają w czasie wielomianowym, podczas gdy czas wersji naiwnej (widoczny tylko dla małych $n$) rośnie wykładniczo.}
	\end{figure}
	
	% ==========================================
	% 2. LCS
	% ==========================================
	\newpage
	\subsection{Najdłuższy Wspólny Podciąg (LCS)}
	Problem LCS polega na znalezieniu najdłuższego ciągu znaków występującego w obu porównywanych napisach w tej samej kolejności (niekoniecznie obok siebie).
	
	\subsubsection{Implementacja}
	Porównano podejście rekurencyjne ze spamiętywaniem oraz iteracyjne.
	\begin{lstlisting}[caption={Fragment implementacji: LCS (wersja iteracyjna)}]
		// Algorytm Bottom-Up o zlozonosci O(m*n).
		// Wypelnia tablice dp, gdzie dp[i][j] to dlugosc LCS dla X[0..i-1] i Y[0..j-1].
		int run_iterative() {
			// Inicjalizacja zerami
			dp_table.assign(m + 1, vector<int>(n + 1, 0));
			
			for (int i = 1; i <= m; ++i) {
				for (int j = 1; j <= n; ++j) {
					// Jesli znaki sa takie same, wydluzamy podciag znaleziony wczesniej (po skosie)
					if (X[i - 1] == Y[j - 1]) {
						dp_table[i][j] = 1 + dp_table[i - 1][j - 1];
					} 
					// W przeciwnym razie bierzemy lepszy wynik z sasiadow (gora lub lewo)
					else {
						dp_table[i][j] = max(dp_table[i - 1][j], dp_table[i][j - 1]);
					}
				}
			}
			table_ready = true;
			return dp_table[m][n];
		}
	\end{lstlisting}
	
	\subsubsection{Wyniki pomiarów}
	Testy przeprowadzono na losowych ciągach znaków (DNA). Jednostka czasu: milisekundy.
	
	% Tabela LCS (Tu mowiles ze dziala, wiec zostawiam jak bylo)
	\begin{table}[H]
		\centering
		\caption{Porównanie czasów dla LCS}
		\pgfplotstabletypeset[
		col sep=comma,
		string type,
		columns={N,Memoized,Iteracyjny},
		columns/N/.style={column name=$n$, column type={r}},
		columns/Memo/.style={column name={Memoized [ms]}, column type={r}},
		columns/Iter/.style={column name={Iteracyjny [ms]}, column type={r}},
		every head row/.style={before row=\toprule,after row=\midrule},
		every last row/.style={after row=\bottomrule},
		]{lcs.csv}
	\end{table}
	
	\begin{figure}[H]
		\centering
		\begin{tikzpicture}
			\begin{axis}[myaxis, title={Wydajność LCS ($O(n^2)$)}, xlabel={$n$}, ylabel={czas [ms]}]
				\addplot+[line width=1pt, color=cPurple, mark=triangle*, mark size=2pt]
				table[x=N, y=Memoized, col sep=comma]{lcs.csv};
				\addlegendentry{Memoized}
				
				\addplot+[line width=1pt, color=cGreen, mark=diamond*, mark size=2pt]
				table[x=N, y=Iteracyjny, col sep=comma]{lcs.csv};
				\addlegendentry{Iteracyjny}
			\end{axis}
		\end{tikzpicture}
		\caption{Wykres potwierdza kwadratową złożoność obliczeniową.}
	\end{figure}
	
	% ==========================================
	% 3. ACTIVITY SELECTOR
	% ==========================================
	\subsection{Wybór zajęć (Activity Selector)}
	Problem polega na wybraniu maksymalnego podzbioru niekonfliktowych zajęć.
	
	\subsubsection{Implementacja}
	Porównano:
	\begin{itemize}
		\item \textbf{Greedy (Zachłanny):} Złożoność zdominowana przez sortowanie $O(n \log n)$. Testowano wersję standardową (sortowanie po czasie zakończenia) i zmodyfikowaną (sortowanie po czasie startu).
		\item \textbf{DP (Dynamiczne):} Złożoność $O(n^2)$, sprawdza wszystkie możliwe podzbiory w sposób usystematyzowany.
	\end{itemize}
	\newpage
	\begin{lstlisting}[caption={Zmodyfikowany greedy}]
		// Algorytm zachlanny dla danych posortowanych po czasie startu.
		// Zlozonosc: O(N log N) (przez sortowanie), sama selekcja O(N).
		void solve_greedy_start_sorted(vector<Activity> acts) {
			// 1. Sortowanie rosnaco po czasie ROZPOCZECIA
			sort(acts.begin(), acts.end(), [](const Activity& a, const Activity& b) {
				return a.start < b.start;
			});
			
			if (acts.empty()) return;
			int n = acts.size();
			
			// Strategia: Wybieramy zajecie o najpozniejszym starcie (ostatnie w posortowanym wektorze)
			// i idziemy "w tyl", szukajac kompatybilnych poprzednikow.
			vector<Activity> result;
			result.push_back(acts[n - 1]); // Pierwszy wybor (od konca)
			int last_added_index = n - 1;
			
			// Iteracja wsteczna
			for (int i = n - 2; i >= 0; --i) {
				// Warunek: koniec aktualnego musi byc <= start poprzednio wybranego
				if (acts[i].finish <= acts[last_added_index].start) {
					result.push_back(acts[i]);
					last_added_index = i;
				}
			}
			// Wynik jest poprawny co do ilosci i zbioru, ale w odwrotnej kolejnosci chronologicznej
		}
	\end{lstlisting}
	
	\subsubsection{Wyniki pomiarów}
	Dla dużych $n$ spodziewana jest drastyczna różnica między podejściem $O(n \log n)$ a $O(n^2)$.
	
	% Tabela Activity - POPRAWIONE KOLUMNY
	\begin{table}[H]
		\centering
		\caption{Porównanie algorytmów Greedy i DP}
		\pgfplotstabletypeset[
		col sep=comma,
		string type,
		% Kolumny z CSV: N, Recursive, Iterative, Start_Sorted, DP
		% Wybieramy do tabeli Iterative i DP (najwazniejsze porownanie)
		columns={N,Iterative,DP},
		columns/N/.style={column name=$n$, column type={r}},
		columns/Iterative/.style={column name={Greedy Iter [ms]}, column type={r}},
		columns/DP/.style={column name={DP [ms]}, column type={r}},
		every head row/.style={before row=\toprule,after row=\midrule},
		every last row/.style={after row=\bottomrule},
		]{activity.csv}
	\end{table}
	
	\begin{figure}[H]
		\centering
		\begin{tikzpicture}
			\begin{axis}[myaxis, title={Greedy vs DP (Activity Selector)}, xlabel={$n$}, ylabel={czas [ms]}]
				
				% Wybieramy Iterative jako reprezentanta Greedy
				\addplot+[line width=1pt, color=cBlue, mark=*, mark size=2pt]
				table[x=N, y=Iterative, col sep=comma]{activity.csv};
				\addlegendentry{Greedy ($O(n \log n)$)}
				
				\addplot+[line width=1pt, color=cPink, mark=x, mark size=2pt]
				table[x=N, y=DP, col sep=comma]{activity.csv};
				\addlegendentry{DP ($O(n^2)$)}
				
			\end{axis}
		\end{tikzpicture}
		\caption{Podejście zachłanne deklasuje programowanie dynamiczne dla tego problemu. Krzywa DP rośnie parabolicznie, podczas gdy Greedy pozostaje blisko osi poziomej.}
	\end{figure}
	
	% ==========================================
	% 4. HUFFMAN
	% ==========================================
	\subsection{Kodowanie Huffmana}
	Zaimplementowano algorytm Huffmana w dwóch wariantach:
	\begin{enumerate}
		\item \textbf{Binarny:} Standardowy, używa alfabetu $\{0, 1\}$.
		\item \textbf{Ternarny:} Zmodyfikowany, używa alfabetu $\{0, 1, 2\}$. Wymaga, aby liczba liści $n$ spełniała warunek $(n-1) \pmod 2 \equiv 0$, w przeciwnym razie dodawane są węzły "dummy".
	\end{enumerate}
	
	\subsubsection{Analiza efektywności kompresji na podstawie danych eksperymentalnych}
	Dla tekstu wejściowego "Lorem ipsum..." (491 znaków) wygenerowano drzewa Huffmana w dwóch wariantach: binarnym i ternarnym. Poniższa tabela przedstawia zestawienie wygenerowanych kodów dla wybranych znaków. 
	
	Kody ternarne są systematycznie krótsze od binarnych. Na przykład dla spacji kod skraca się z 3 bitów do 2 trytów, a dla rzadkiego znaku 'Q' z 9 bitów do 5 trytów.
	
	\begin{table}[H]
		\centering
		\caption{Bezpośrednie porównanie kodów binarnych i ternarnych (wybrane znaki)}
		\begin{tabular}{lllc}
			\toprule
			\textbf{Znak} & \textbf{Kod Binarny (0-1)} & \textbf{Kod Ternarny (0-2)} & \textbf{Zysk długości} \\
			\midrule
			' ' (spacja) & \texttt{101} (3) & \texttt{20} (2) & -1 \\
			'e' & \texttt{001} (3) & \texttt{11} (2) & -1 \\
			'a' & \texttt{0101} (4) & \texttt{212} (3) & -1 \\
			's' & \texttt{1100} (4) & \texttt{00} (2) & -2 \\
			'p' & \texttt{111101} (6) & \texttt{2222} (4) & -2 \\
			'L' & \texttt{100111010} (9) & \texttt{02101} (5) & -4 \\
			'Q' & \texttt{100111011} (9) & \texttt{02102} (5) & -4 \\
			\bottomrule
		\end{tabular}
	\end{table}
	
	\noindent
	Całkowite długości zakodowanych ciągów dla całego tekstu przedstawiają się następująco:
	
	\begin{table}[H]
		\centering
		\caption{Sumaryczna długość zakodowanego tekstu}
		\begin{tabular}{lr}
			\toprule
			\textbf{Metoda} & \textbf{Całkowita długość} \\
			\midrule
			Huffman Binarny & 2058 bitów \\
			Huffman Ternarny & \textbf{1311 trytów} \\
			\bottomrule
		\end{tabular}
	\end{table}
	
	\noindent
	Stosunek długości wyniósł:
	$$ \text{Ratio} = \frac{1311}{2058} \approx 0.637 $$
	
	\noindent
	\textbf{Wnioski:} Wynik eksperymentalny ($0.637$) potwierdza teorię informacji. Jeden tryt przenosi więcej informacji niż bit ($\log_2 3 \approx 1.585$ bita). Teoretycznie idealna kompresja pozwoliłaby skrócić zapis do ok. $63.1\%$ oryginalnej długości wyrażonej w bitach. Uzyskany wynik jest bardzo bliski tej granicy.
	\section{Wnioski końcowe}
	Przeprowadzone eksperymenty potwierdziły, że dobór odpowiedniego algorytmu jest kluczowy dla wydajności systemu.
	\begin{itemize}
		\item W problemach takich jak \textit{Rod Cutting} czy \textit{LCS}, zastosowanie programowania dynamicznego pozwala rozwiązywać problemy w czasie wielomianowym, które naiwnie wymagałyby czasu wykładniczego.
		\item W przypadku \textit{Activity Selector}, istnienie własności wyboru zachłannego pozwala zredukować złożoność do $O(n \log n)$, co czyni użycie DP ($O(n^2)$) nieefektywnym.
		\item Dla kodowania Huffmana, zwiększenie podstawy systemu kodowania skraca długość wynikowego napisu, co może być korzystne w przesyłaniu danych.
	\end{itemize}
	
\end{document}